% Template for PLoS
% Version 3.5 March 2018
%
% % % % % % % % % % % % % % % % % % % % % %
%
% -- IMPORTANT NOTE
%
% This template contains comments intended
% to minimize problems and delays during our production
% process. Please follow the template instructions
% whenever possible.
%
% % % % % % % % % % % % % % % % % % % % % % %
%
% Once your paper is accepted for publication,
% PLEASE REMOVE ALL TRACKED CHANGES in this file
% and leave only the final text of your manuscript.
% PLOS recommends the use of latexdiff to track changes during review, as this will help to maintain a clean tex file.
% Visit https://www.ctan.org/pkg/latexdiff?lang=en for info or contact us at latex@plos.org.
%
%
% There are no restrictions on package use within the LaTeX files except that
% no packages listed in the template may be deleted.
%
% Please do not include colors or graphics in the text.
%
% The manuscript LaTeX source should be contained within a single file (do not use \input, \externaldocument, or similar commands).
%
% % % % % % % % % % % % % % % % % % % % % % %
%
% -- FIGURES AND TABLES
%
% Please include tables/figure captions directly after the paragraph where they are first cited in the text.
%
% DO NOT INCLUDE GRAPHICS IN YOUR MANUSCRIPT
% - Figures should be uploaded separately from your manuscript file.
% - Figures generated using LaTeX should be extracted and removed from the PDF before submission.
% - Figures containing multiple panels/subfigures must be combined into one image file before submission.
% For figure citations, please use "Fig" instead of "Figure".
% See http://journals.plos.org/plosone/s/figures for PLOS figure guidelines.
%
% Tables should be cell-based and may not contain:
% - spacing/line breaks within cells to alter layout or alignment
% - do not nest tabular environments (no tabular environments within tabular environments)
% - no graphics or colored text (cell background color/shading OK)
% See http://journals.plos.org/plosone/s/tables for table guidelines.
%
% For tables that exceed the width of the text column, use the adjustwidth environment as illustrated in the example table in text below.
%
% % % % % % % % % % % % % % % % % % % % % % % %
%
% -- EQUATIONS, MATH SYMBOLS, SUBSCRIPTS, AND SUPERSCRIPTS
%
% IMPORTANT
% Below are a few tips to help format your equations and other special characters according to our specifications. For more tips to help reduce the possibility of formatting errors during conversion, please see our LaTeX guidelines at http://journals.plos.org/plosone/s/latex
%
% For inline equations, please be sure to include all portions of an equation in the math environment.
%
% Do not include text that is not math in the math environment.
%
% Please add line breaks to long display equations when possible in order to fit size of the column.
%
% For inline equations, please do not include punctuation (commas, etc) within the math environment unless this is part of the equation.
%
% When adding superscript or subscripts outside of brackets/braces, please group using {}.
%
% Do not use \cal for caligraphic font.  Instead, use \mathcal{}
%
% % % % % % % % % % % % % % % % % % % % % % % %
%
% Please contact latex@plos.org with any questions.
%
% % % % % % % % % % % % % % % % % % % % % % % %

\documentclass[10pt,letterpaper]{article}
\usepackage[top=0.85in,left=2.75in,footskip=0.75in]{geometry}

% amsmath and amssymb packages, useful for mathematical formulas and symbols
\usepackage{amsmath,amssymb}

% Use adjustwidth environment to exceed column width (see example table in text)
\usepackage{changepage}

% textcomp package and marvosym package for additional characters
\usepackage{textcomp,marvosym}

% cite package, to clean up citations in the main text. Do not remove.
\usepackage{cite}

% Use nameref to cite supporting information files (see Supporting Information section for more info)
\usepackage{nameref,hyperref}

% line numbers
\usepackage[right]{lineno}

% ligatures disabled
\usepackage{microtype}
\DisableLigatures[f]{encoding = *, family = * }

% color can be used to apply background shading to table cells only
\usepackage[table]{xcolor}

% array package and thick rules for tables
\usepackage{array}

% create "+" rule type for thick vertical lines
\newcolumntype{+}{!{\vrule width 2pt}}

% create \thickcline for thick horizontal lines of variable length
\newlength\savedwidth
\newcommand\thickcline[1]{%
  \noalign{\global\savedwidth\arrayrulewidth\global\arrayrulewidth 2pt}%
  \cline{#1}%
  \noalign{\vskip\arrayrulewidth}%
  \noalign{\global\arrayrulewidth\savedwidth}%
}

% \thickhline command for thick horizontal lines that span the table
\newcommand\thickhline{\noalign{\global\savedwidth\arrayrulewidth\global\arrayrulewidth 2pt}%
\hline
\noalign{\global\arrayrulewidth\savedwidth}}


% Remove comment for double spacing
%\usepackage{setspace}
%\doublespacing

% Text layout
\raggedright
\setlength{\parindent}{0.5cm}
\textwidth 5.25in
\textheight 8.75in

% Bold the 'Figure #' in the caption and separate it from the title/caption with a period
% Captions will be left justified
\usepackage[aboveskip=1pt,labelfont=bf,labelsep=period,justification=raggedright,singlelinecheck=off]{caption}
\renewcommand{\figurename}{Fig}

% Use the PLoS provided BiBTeX style
\bibliographystyle{plos2015}

% Remove brackets from numbering in List of References
\makeatletter
\renewcommand{\@biblabel}[1]{\quad#1.}
\makeatother



% Header and Footer with logo
\usepackage{lastpage,fancyhdr,graphicx}
\usepackage{epstopdf}
%\pagestyle{myheadings}
\pagestyle{fancy}
\fancyhf{}
%\setlength{\headheight}{27.023pt}
%\lhead{\includegraphics[width=2.0in]{PLOS-submission.eps}}
\rfoot{\thepage/\pageref{LastPage}}
\renewcommand{\headrulewidth}{0pt}
\renewcommand{\footrule}{\hrule height 2pt \vspace{2mm}}
\fancyheadoffset[L]{2.25in}
\fancyfootoffset[L]{2.25in}
\lfoot{\today}

%% Include all macros below

% ------------------------------------------------------------------
% Added for this manuscript (adding packages is allowed by PLOS).
% float: the [H] placement keeps every table and figure caption
% directly after the paragraph where it is first cited, as PLOS asks.
\usepackage{float}

% PLOS wants figures uploaded as separate TIFF files, NOT inside the PDF.
% Keep \showfigsfalse for the submission PDF.
% Change to \showfigstrue only to make a private reading copy with the figures
% (they are read from ../figures/, the folder made by the Python scripts).
\newif\ifshowfigs
\showfigsfalse
\graphicspath{{../figures/}{figures/}}
\newcommand{\plosfig}[1]{\ifshowfigs\par\includegraphics[width=\linewidth]{#1}\par\fi}
% ------------------------------------------------------------------

\begin{document}
\vspace*{0.2in}

% Title must be 250 characters or less.
\begin{flushleft}
{\Large
\textbf\newline{Fitted survival models and the molecular international prognostic scoring system in myeloid neoplasms: a centre-grouped validation study with an unseen external centre} % Please use "sentence case" for title and headings (capitalize only the first word in a title (or heading), the first word in a subtitle (or subheading), and any proper nouns).
}
\newline
% Insert author names, affiliations and corresponding author email (do not include titles, positions, or degrees).
\\
Bernard Frimpong\textsuperscript{1*}
\\
\bigskip
\textbf{1} Independent researcher, Akron, Ohio, United States of America
\\
\bigskip

% Use the asterisk to denote corresponding authorship and provide email address in note below.
* bernardfrimpong040@gmail.com

\end{flushleft}
% Please keep the abstract below 300 words
\section*{Abstract}
\textbf{Background:} Machine-learning survival models are increasingly proposed for
prognosis in myeloid neoplasms, but they are rarely compared with the current molecular
standard, the Molecular International Prognostic Scoring System (IPSS-M), under a
validation design that mirrors use at a new hospital.

\textbf{Methods:} I analysed 3,173 patients from 23 centres (1,600 deaths) released in a
public prediction challenge, and an external centre of 1,193 patients whose outcomes were
withheld. Predictors were blood counts, a rule-based reading of free-text karyotypes, and
somatic mutations restricted to the 64 genes sequenced at both the training and the
external centres. The published revised score (IPSS-R) and the IPSS-M, applied without
refitting, were compared with Cox models, an elastic-net Cox model, a random survival
forest, gradient-boosted Cox trees and a rank-averaged ensemble. Discrimination was
measured by the censoring-weighted concordance index truncated at 7 years, over 15
centre-grouped cross-validation folds with nested tuning and corrected confidence
intervals, and on the external centre.

\textbf{Results:} Without fitting, the IPSS-M reached a concordance of 0.720 (95\%
confidence interval 0.693 to 0.748), against 0.669 for the IPSS-R. The best fitted model,
an ensemble of the elastic-net and gradient-boosted Cox models chosen before external
scoring, reached 0.725 (0.700 to 0.750) and beat the elastic-net alone on all 15 folds,
but its advantage over the IPSS-M was small and uncertain (+0.005, $-$0.005 to +0.014;
p = 0.29). A recalibrated IPSS-M was slightly better calibrated (integrated Brier score
0.160 against 0.163). On the external centre, the ensemble scored 0.755 and the IPSS-M
0.726, a gap unexplained by case mix or missing data. Mutation burden on the shared gene panel was the most important
predictor group.

\textbf{Conclusions:} Fitted models matched the IPSS-M internally and exceeded it at one
external centre. New prognostic models for myeloid neoplasms should be benchmarked
against the IPSS-M with centre-level validation before improvement is claimed.

% Please keep the Author Summary between 150 and 200 words
% Use first person. PLOS ONE authors please skip this step.
% Author Summary not valid for PLOS ONE submissions.

\linenumbers

% Use "Eq" instead of "Equation" for equation citations.
\section*{Introduction}
Myeloid neoplasms, mainly the myelodysplastic syndromes (MDS) and acute myeloid leukemia,
range from slowly progressing disease to rapidly fatal illness. Estimating a patient's
prognosis at diagnosis guides treatment: patients at low risk may receive supportive care,
while patients at high risk may be considered for intensive therapy such as allogeneic
stem cell transplantation. For MDS, risk has long been estimated with the revised
International Prognostic Scoring System (IPSS-R), which combines the karyotype, bone-marrow
blasts and three blood counts~\cite{greenberg2012,schanz2012}. Its molecular successor, the
IPSS-M, adds somatic mutations in 31 genes and was developed and validated on several
thousand patients~\cite{bernard2022}.

At the same time, machine-learning survival models such as random survival
forests~\cite{ishwaran2008} and gradient-boosted trees~\cite{friedman2001,chen2016} are
often proposed as improvements over regression. Two features of the evidence call for
caution. First, systematic reviews of clinical prediction models have found no average
performance benefit of machine learning over regression when both are developed
carefully~\cite{christodoulou2019}. Second, new models are frequently compared with weak
baselines and validated on random splits of a single dataset, which mixes patients from
the same hospital into the training and validation sets; such internal validation
overstates performance at a new hospital, where case mix and measurement practice
differ~\cite{austin2016,steyerberg2019}.

A public data challenge organised by Qube Research and Technologies (QRT) and Institut
Gustave Roussy offers a setting to test both points~\cite{challenge162}. It provides
clinical, cytogenetic and molecular data on patients from 23 centres, and scores
predictions on a further centre whose outcomes are never released. A similar crowdsourced
effort, the DREAM 9 acute myeloid leukemia challenge, also scored survival models by a
concordance index~\cite{noren2016}.

The aim of this study was to ask whether fitted survival models, from penalised Cox
regression to boosted trees and their ensemble, improve on the published IPSS-M when
validated on hospitals they were not trained on. The study is reported following the
TRIPOD+AI guideline~\cite{collins2024} (\nameref{S1_Table}).

\section*{Materials and methods}

\subsection*{Data source and participants}
The data are those of challenge 162 on the Challenge Data platform (\'Ecole normale
sup\'erieure, Paris)~\cite{challenge162}. The training set contains 3,323 patients from 23
clinical centres, with overall survival in years and a death indicator. The test set
contains 1,193 patients, all from one further centre (coded KYW) that contributes no
training patients; their outcomes are withheld, and predictions for them are scored by the
platform. One hundred and fifty training patients have no recorded outcome and were
excluded, leaving 3,173 patients and 1,600 deaths. The median blast count of 3\% indicates
that most patients have MDS. Age, sex, disease subtype and treatment are not provided.

\subsection*{Ethics statement}
This is a secondary analysis of de-identified data distributed by the challenge organisers
to registered participants. No new data were collected, and the author had no access to
identifying information. [TO CONFIRM WITH THE ORGANISERS: the name of the ethics committee
that approved the original data collection, whether patient consent was written, and
written permission to publish results derived from the challenge data.]

\subsection*{Outcome}
The outcome was overall survival, the time from diagnosis to death from any cause, with
patients alive at last follow-up treated as right-censored. Fifty-three patients had an
observed time of zero; all were censored and left the risk set immediately.

\subsection*{Predictors}
Ninety-one candidate predictors were built from each patient's own data; full definitions
are given in \nameref{S1_Appendix}.

\paragraph*{Clinical variables.} Bone-marrow blasts (\%), haemoglobin (g/dL), and the
logarithm of one plus white cells, neutrophils, monocytes and platelets (G/L), with a 0/1
flag for each missing value. Missing values were replaced by the median of the training
part of each fold, inside the model pipeline.

\paragraph*{Karyotype.} Karyotypes are free text in International System for Human
Cytogenomic Nomenclature notation, written differently by different centres (lower case
in training, upper case and the word ``Normal'' at the test centre). A rule-based parser
normalised the text and extracted the number of abnormalities, 14 recurrent lesions
(including monosomy 7, deletion 5q, trisomy 8 and 17p loss), complex and monosomal status,
and the IPSS-R cytogenetic category~\cite{schanz2012}. It was checked on 25 hand-written
test strings.

\paragraph*{Mutations.} The training file lists 124 mutated genes and the test file 83, of
which only 64 appear in both; several genes frequent in training never appear at the test
centre, which suggests a smaller sequencing panel there. To give mutation counts the same
meaning at every centre, all molecular features used only the 64 shared genes. Features were
the numbers of mutations and of mutated genes, the maximum variant allele fraction (VAF),
the number of truncating mutations, TP53 count, maximum VAF and multi-hit status, FLT3
internal tandem duplication, MLL partial tandem duplication, six pathway counts, and
indicators for the 45 shared genes mutated in at least 1\% of training patients.

\subsection*{Benchmark scores}
Both benchmark scores were computed from their published definitions, with no parameter
estimated from these data.

\paragraph*{IPSS-R.} The sum of points for cytogenetic category, blasts, haemoglobin,
platelets and neutrophils~\cite{greenberg2012}; a missing component received the median
points of the training patients.

\paragraph*{IPSS-M.} The IPSS-M is a weighted sum of centred variables,
\begin{equation}
\text{IPSS-M} = \frac{1}{\ln 2}\sum_{k} \beta_k \left(x_k - \mu_k\right),
\label{eq:ipssm}
\end{equation}
where $\beta_k$ is the published coefficient of variable $k$ and $\mu_k$ its published
reference mean~\cite{bernard2022}. The coefficients and means were taken from the authors'
open-source calculator~\cite{ipssmcalc}. A missing input was set to its reference mean,
so that it contributes nothing to the sum. The implementation agreed exactly with the
official calculator on 300 training patients with complete data. Copy-neutral loss of
heterozygosity is not available in these data, so a single TP53 mutation counted as
multi-hit only when its VAF exceeded 0.55 or a 17p loss was present.

\subsection*{Models}
All fitted models used the same predictors and were fitted with the same procedure in every
fold.

\paragraph*{Cox models.} The Cox proportional hazards model~\cite{cox1972} writes the hazard
of a patient with predictors $x$ as
\begin{equation}
h(t \mid x) = h_0(t)\exp\left(\beta^{\top} x\right),
\label{eq:cox}
\end{equation}
and the linear predictor $\beta^{\top} x$ is the risk score. Three versions were fitted: a
four-variable model (blasts, haemoglobin, platelets and number of mutations) reconstructed
from the challenge description; an IPSS-M-style model with 24 variables chosen from the
literature (the five blood values, the IPSS-R cytogenetic category, a missing-karyotype
flag, the 16 IPSS-M main-effect genes and the capped count of IPSS-M residual genes), with
a small ridge penalty of 0.1; and an elastic-net Cox model~\cite{zou2005,simon2011} on 84
predictors (all except monocytes and the missing flags), which maximises
\begin{equation}
\ln L(\beta) - \alpha\left(\rho \sum_{k} |\beta_k| + \frac{1-\rho}{2}\sum_{k} \beta_k^{2}\right),
\label{eq:enet}
\end{equation}
where $L$ is the Cox partial likelihood. The penalty strength $\alpha$ (30 values) and mix
$\rho \in \{0.1, 0.5, 1\}$ were tuned by nested cross-validation.

\paragraph*{Tree models.} A random survival forest~\cite{ishwaran2008} with 300 trees,
at least 15 patients per leaf and nine predictors tried per split; and gradient-boosted
trees with a Cox loss (XGBoost~\cite{chen2016}; learning rate 0.03, 80\% of patients and
50\% of predictors sampled per tree, minimum child weight 10). The boosted model's depth
(2, 3 or 4) and number of trees (50 to 600) were tuned by nested cross-validation, and five
models with different random seeds were averaged. Exact split finding on a single thread
was used, so that results are reproducible on one machine.

\paragraph*{Ensemble.} Because the metric depends only on the order of the risk scores,
each model's scores were converted to ranks between 0 and 1 and averaged. The ensemble of
the elastic-net and boosted Cox models was specified as the final model before any of its
external scores were known.

\subsection*{Validation}
\paragraph*{Centre-grouped cross-validation.} All patients of a centre were placed in the
same fold, so that every validation fold contained only centres unseen in training. The 23
centres were dealt into five folds of similar size in a random order, and this was repeated
three times with different orders, giving 15 validation folds. All tuning was nested: the
penalty and tree settings were chosen by an inner four-fold split of the training part,
also grouped by centre, and median imputation and standardisation were learned on the
training part only. Every model was evaluated on the same 15 folds, so that models could be
compared fold by fold.

\paragraph*{External validation.} Models refitted on all 3,173 training patients were
applied to the test centre and scored by the challenge platform on its public subset of
patients. Submissions are limited to two per day. Before each of the last submissions, the
expected score was written down: the public score of the elastic-net plus the model's
cross-validated paired difference with the elastic-net.

\subsection*{Performance measures}
Discrimination was measured by the inverse-probability-of-censoring-weighted concordance
index~\cite{uno2011}, the challenge metric. For patients $i$ with observed time $T_i$,
death indicator $\Delta_i$ and risk score $\eta_i$,
\begin{equation}
C_{\tau} = \frac{\sum_{i}\sum_{j} \Delta_i\,\hat{G}(T_i)^{-2}\,
\mathbf{1}(T_i < T_j,\ T_i < \tau)\,\mathbf{1}(\eta_i > \eta_j)}
{\sum_{i}\sum_{j} \Delta_i\,\hat{G}(T_i)^{-2}\,\mathbf{1}(T_i < T_j,\ T_i < \tau)},
\label{eq:uno}
\end{equation}
where $\hat{G}$ is the Kaplan--Meier estimate~\cite{kaplan1958} of the probability of
remaining uncensored, estimated on the training part, and $\tau = 7$ years. A pair started
by a death at a time when few patients remain under follow-up is rare to observe and so
receives a larger weight. Harrell's concordance index~\cite{harrell1982} was also recorded.
The cumulative/dynamic area under the receiver operating characteristic curve at 1, 3 and 5
years~\cite{hung2010} was computed on all 15 folds, with validation follow-up censored at
6.5 years so that the censoring weights stay defined. Calibration was assessed on the five
folds of the first repetition by the Brier score at 3 years and integrated from 0.25 to 6
years~\cite{graf1999}, and by comparing mean predicted with Kaplan--Meier observed 3-year
survival in ten groups of predicted risk~\cite{vancalster2019}. For this comparison the
IPSS-M was recalibrated by a one-variable Cox model on its score, fitted on each training
part, so that its published weights were kept and only the link from score to probability
was estimated; the elastic-net used its Breslow baseline survival.

\subsection*{Statistical analysis}
The 15 cross-validation scores are not independent, because the three repetitions reuse the
same patients and training sets overlap. Confidence intervals and two-sided p-values for
mean scores and paired differences therefore used the corrected resampled t-test of Nadeau
and Bengio~\cite{nadeau2003}, which inflates the variance of the mean to
\begin{equation}
\widehat{\operatorname{Var}} = \left(\frac{1}{J} + \frac{n_{\text{valid}}}{n_{\text{train}}}\right) s^{2}
= \left(\frac{1}{15} + 0.257\right) s^{2} = 0.323\, s^{2},
\label{eq:nb}
\end{equation}
where $J = 15$ is the number of folds, $s$ the standard deviation of the fold scores, and
0.257 the average ratio of validation to training patients, with a t distribution on 14
degrees of freedom. The proportional hazards assumption of the IPSS-M-style Cox model was
checked with scaled Schoenfeld residuals~\cite{grambsch1994}. Permutation
importance~\cite{breiman2001} was measured on held-out folds by shuffling groups of related
predictors together and recording the fall in concordance.

To interpret the difference between the ensemble and the IPSS-M at the external centre,
three explanations were tested on the training data. For case mix, each validation fold was
resampled 20 times without replacement to the external centre's distribution of IPSS-R risk
groups. For missing data, blood counts and blast percentages were removed for an extra 7\%
of each validation fold, matching the higher missingness at the external centre, and the
fall in concordance of the IPSS-M and the elastic-net was compared. For between-centre
variation, the difference was computed within each training centre with at least 100
patients, using out-of-fold predictions.

\subsection*{Software}
Analyses used Python 3.13 with pandas [VERSION], NumPy [VERSION], scikit-learn
[VERSION]~\cite{pedregosa2011}, scikit-survival [VERSION]~\cite{polsterl2020},
lifelines [VERSION]~\cite{lifelines2019} and XGBoost 3.4.1~\cite{chen2016}. All code is
available at https://github.com/bernardfrimpong040-pixel/aml-survival-prediction
[ARCHIVE ON ZENODO AND ADD THE DOI BEFORE SUBMISSION].

\subsection*{Use of artificial intelligence tools}
An AI assistant (Claude, Anthropic) was used to help write and review analysis code and
to edit the manuscript text. The author designed the study, ran every analysis, checked
all results and takes full responsibility for the content.

\section*{Results}

\subsection*{Cohort}
Table~\ref{table1} describes the two cohorts. The median overall survival in training was
3.15 years. The test centre was clinically sicker than the training centres: it had more
blasts, lower haemoglobin and platelets, almost three times the share of complex
karyotypes, and nearly three times the share of TP53 multi-hit patients. One third of its
patients were IPSS-R very high risk, against 11\% in training. Monocytes were missing for
74\% of test patients.

\begin{table}[H]
\begin{adjustwidth}{-2.25in}{0in}
\centering
\caption{{\bf Characteristics of the training cohort and the external test centre.}}
\begin{tabular}{|l|l|l|}
\hline
{\bf Characteristic} & {\bf Training (23 centres)} & {\bf Test (one centre)}\\ \thickhline
Patients & 3,173 & 1,193 \\ \hline
Sex inferred from karyotype: female / male / unknown, \% & 33.3 / 52.7 / 14.0 & 17.1 / 45.8 / 37.1 \\ \hline
Deaths, n (\%) & 1,600 (50.4) & withheld \\ \hline
Median overall survival, years & 3.15 & withheld \\ \hline
Median follow-up (reverse Kaplan--Meier), years & 3.84 & withheld \\ \hline
Bone-marrow blasts, \% & 3.0 (1.0--8.0) & 6.0 (2.0--17.0) \\ \hline
Haemoglobin, g/dL & 9.7 (8.5--11.1) & 8.4 (7.4--10.1) \\ \hline
Platelets, G/L & 123 (65--230) & 74 (37--143) \\ \hline
White blood cells, G/L & 4.1 (2.7--6.7) & 2.9 (1.9--5.3) \\ \hline
Neutrophils, G/L & 2.00 (1.00--3.70) & 1.30 (0.54--2.64) \\ \hline
Karyotype missing or unusable, n (\%) & 350 (11.0) & 120 (10.1) \\ \hline
Normal karyotype, n (\%) & 1,643 (51.8) & 410 (34.4) \\ \hline
Complex karyotype, n (\%) & 307 (9.7) & 336 (28.2) \\ \hline
IPSS-R cytogenetics poor or very poor, n (\%) & 408 (12.9) & 380 (31.9) \\ \hline
Mutated genes per patient (shared panel) & 2 (1--4) & 2 (1--3) \\ \hline
TP53 multi-hit, n (\%) & 210 (6.6) & 213 (17.9) \\ \hline
SF3B1 mutated, n (\%) & 727 (22.9) & 92 (7.7) \\ \hline
IPSS-R very high risk, n (\%) & 363 (11.4) & 395 (33.1) \\ \hline
IPSS-M very high risk, n (\%) & 666 (21.0) & 494 (41.4) \\ \hline
\end{tabular}
\begin{flushleft} Values are median (first--third quartile) unless stated. Sex is not
recorded in the data and was inferred from the sex chromosomes in the karyotype. Outcomes
of the test centre are withheld by the challenge.
\end{flushleft}
\label{table1}
\end{adjustwidth}
\end{table}

Both benchmark scores separated survival in the expected order. The six IPSS-M categories
gave a log-rank chi-squared statistic of 974.2 on 5 degrees of freedom, against 737.4 on 4
for the five IPSS-R groups, with median survival falling from 9.65 years in the IPSS-M very
low category to 1.03 years in the very high category (Fig~\ref{fig1}). At the test centre,
41\% of patients were IPSS-M very high risk, against 21\% in training.

\begin{figure}[H]
\plosfig{fig25_ipssm.png}
\caption{{\bf Survival by IPSS-M category and case mix of the two cohorts.}
Left: Kaplan--Meier overall survival by IPSS-M category in the 3,173 training patients.
Right: share of patients in each IPSS-M category in training and at the external test
centre.}
\label{fig1}
\end{figure}

\subsection*{Discrimination in centre-grouped cross-validation}
Table~\ref{table2} gives the concordance of every model over the 15 centre-grouped folds.
The IPSS-R reached 0.669, and its spread across folds (standard deviation 0.029, range
0.624 to 0.729), which involves no fitting, shows how much the difficulty of the task
varies between centres. The IPSS-M, also without fitting, reached 0.720. Of the fitted
models, the four-variable Cox model reached 0.688 and the IPSS-M-style Cox model 0.716.
The elastic-net reached 0.722, the random survival forest 0.719, the boosted Cox model
0.724 and the ensemble 0.725. Adding monocytes and the missing-value flags to the
elastic-net changed its concordance by $-0.0004$ and lost 11 of 15 folds, so these
variables were left out.

\begin{table}[H]
\begin{adjustwidth}{-2.25in}{0in}
\centering
\caption{{\bf Concordance over 15 centre-grouped cross-validation folds.}}
\begin{tabular}{|l|c|c|c|c|}
\hline
{\bf Model} & {\bf Fitted} & {\bf IPCW C-index (95\% CI)} & {\bf SD over folds} & {\bf Harrell's C}\\ \thickhline
IPSS-R score & No & 0.669 (0.634--0.703) & 0.029 & 0.692 \\ \hline
IPSS-M score & No & 0.720 (0.693--0.748) & 0.023 & -- \\ \hline
Four-variable Cox & Yes & 0.688 (0.653--0.723) & 0.029 & 0.703 \\ \hline
IPSS-M-style Cox (24 variables) & Yes & 0.716 (0.688--0.743) & 0.023 & 0.736 \\ \hline
Elastic-net Cox (84 predictors) & Yes & 0.722 (0.696--0.747) & 0.021 & 0.740 \\ \hline
Random survival forest & Yes & 0.719 & 0.021 & -- \\ \hline
Boosted Cox trees (XGBoost) & Yes & 0.724 (0.700--0.749) & 0.020 & -- \\ \hline
Ensemble: elastic-net and boosted Cox & Yes & 0.725 (0.700--0.750) & 0.021 & -- \\ \hline
\end{tabular}
\begin{flushleft} IPCW C-index, concordance index weighted by the inverse probability of
censoring, truncated at 7 years. Confidence intervals use the corrected resampled t-test
(Eq~\ref{eq:nb}). SD, standard deviation. Dashes: not computed.
\end{flushleft}
\label{table2}
\end{adjustwidth}
\end{table}

Because all models were scored on the same folds, paired differences remove most of this
between-centre variation (Table~\ref{table3}, Fig~\ref{fig2}). The gains from clinical
knowledge were large and certain: the IPSS-M beat the IPSS-R by 0.052 (95\% CI 0.037 to
0.066), and the 24 literature-chosen variables beat the four-variable model by 0.028 (0.010
to 0.045), winning 14 of 15 folds. The gains from more flexible algorithms were small. The
ensemble beat the elastic-net by 0.003 (0.002 to 0.005) on all 15 folds; the boosted model
alone beat it by 0.003 ($-$0.000 to 0.005) on 14 of 15. No fitted model was distinguishable
from the IPSS-M: the ensemble was ahead by 0.005 ($-$0.005 to 0.014; p = 0.29), winning 11
of 15 folds, and the elastic-net by 0.002 ($-$0.007 to 0.011; p = 0.71).

\begin{table}[H]
\begin{adjustwidth}{-2.25in}{0in}
\centering
\caption{{\bf Paired differences in concordance over the same 15 folds.}}
\begin{tabular}{|l|c|c|c|}
\hline
{\bf Comparison} & {\bf Mean difference (95\% CI)} & {\bf p} & {\bf Folds won}\\ \thickhline
IPSS-M minus IPSS-R & 0.052 (0.037 to 0.066) & $<$0.001 & 15/15 \\ \hline
IPSS-M-style Cox minus four-variable Cox & 0.028 (0.010 to 0.045) & 0.004 & 14/15 \\ \hline
Elastic-net minus IPSS-R & 0.053 (0.033 to 0.073) & $<$0.001 & 15/15 \\ \hline
Elastic-net minus IPSS-M & 0.002 ($-$0.007 to 0.011) & 0.71 & 10/15 \\ \hline
Boosted Cox minus elastic-net & 0.003 ($-$0.000 to 0.005) & 0.07 & 14/15 \\ \hline
Ensemble minus elastic-net & 0.003 (0.002 to 0.005) & $<$0.001 & 15/15 \\ \hline
Ensemble minus IPSS-M & 0.005 ($-$0.005 to 0.014) & 0.29 & 11/15 \\ \hline
\end{tabular}
\begin{flushleft} Positive values favour the first model. Confidence intervals and
two-sided p-values from the corrected resampled t-test with 14 degrees of freedom.
\end{flushleft}
\label{table3}
\end{adjustwidth}
\end{table}

\begin{figure}[H]
\plosfig{fig26_paired_ci.png}
\caption{{\bf Paired differences in concordance with corrected 95\% confidence intervals.}
Each row is the mean, over the same 15 centre-grouped folds, of the difference in the
censoring-weighted concordance index between two models.}
\label{fig2}
\end{figure}

The time-dependent area under the curve told the same story (\nameref{S4_Fig}). At 1, 3
and 5 years, the fitted models and the IPSS-M were within 0.006 of each other (for the
ensemble 0.779, 0.805 and 0.806; for the IPSS-M 0.777, 0.803 and 0.803), and 0.05 to 0.08
above the IPSS-R.

\subsection*{What the models use}
In the IPSS-M-style Cox model refitted on all training patients, TP53 multi-hit status had
the largest hazard ratio (3.21, 95\% CI 2.59 to 3.98), followed by MLL partial tandem
duplication (2.05) and FLT3 mutation (1.65), while higher haemoglobin (0.77 per standard
deviation) and platelets (0.83) were protective (Fig~\ref{fig3}). SF3B1 mutation, strongly
protective in a univariable comparison, was neutral after adjustment (1.02, 0.88 to 1.18).
Four of 24 variables departed from proportional hazards at p $<$ 0.05; the largest
departure was NPM1 mutation, harmful in the first months and protective after about four
years (\nameref{S3_Fig}).

\begin{figure}[H]
\plosfig{fig15_hazard_ratios.png}
\caption{{\bf Hazard ratios of the IPSS-M-style Cox model.}
Hazard ratios with 95\% confidence intervals, refitted on all 3,173 training patients.
Blood values are per standard deviation; the cytogenetic category and the residual-gene
count are per step; genes are mutated against not mutated.}
\label{fig3}
\end{figure}

Permutation importance on held-out centres identified mutation burden (the numbers of
mutations and mutated genes, the maximum VAF and the number of truncating mutations, all
counted on the shared panel) as the most important predictor group for both the
elastic-net and the boosted model, with a fall in concordance of about 0.04 when shuffled,
more than twice that of any other group (Fig~\ref{fig4}). The boosted model depended less
on the TP53 columns than the elastic-net (0.007 against 0.016), consistent with the strong
overlap between TP53 multi-hit status and complex karyotype: 88.5\% of TP53 multi-hit
patients had a complex karyotype, against 4.7\% of the others.

\begin{figure}[H]
\plosfig{fig23_permutation_importance.png}
\caption{{\bf Permutation importance on held-out centres.}
Fall in the censoring-weighted concordance index when the predictors of one group are
shuffled together in the validation fold (mean and standard deviation over the five
folds of the first repetition), for the elastic-net and the boosted Cox model.}
\label{fig4}
\end{figure}

\subsection*{Calibration}
Both the recalibrated IPSS-M and the elastic-net reduced the Brier score by about a quarter
compared with a Kaplan--Meier curve that ignores all predictors (Table~\ref{table4}). The
recalibrated IPSS-M was slightly better calibrated. The elastic-net's predictions were
compressed towards the middle: its predicted 3-year survival spanned 0.08 to 0.74 across
the ten risk groups, while observed survival spanned 0.08 to 0.85, so patients at low risk
did better than predicted (Fig~\ref{fig5}). Such compression is the expected effect of
penalisation and could be corrected by recalibrating the linear predictor.

\begin{table}[H]
\begin{adjustwidth}{-2.25in}{0in}
\centering
\caption{{\bf Brier scores and calibration (five folds of the first repetition).}}
\begin{tabular}{|l|c|c|c|}
\hline
{\bf Model} & {\bf Integrated Brier score} & {\bf Brier score at 3 years} & {\bf Calibration error at 3 years}\\ \thickhline
Kaplan--Meier, no predictors & 0.211 & 0.243 & -- \\ \hline
IPSS-M, recalibrated & 0.160 & 0.178 & 0.050 \\ \hline
Elastic-net Cox & 0.163 & 0.181 & 0.065 \\ \hline
\end{tabular}
\begin{flushleft} Lower is better. The integrated Brier score averages the Brier score from
0.25 to 6 years. Calibration error is the mean absolute difference between predicted and
Kaplan--Meier observed 3-year survival over ten groups of predicted risk.
\end{flushleft}
\label{table4}
\end{adjustwidth}
\end{table}

\begin{figure}[H]
\plosfig{fig28_calibration.png}
\caption{{\bf Calibration at 3 years.}
Mean predicted against Kaplan--Meier observed 3-year survival in ten groups of
out-of-fold predicted risk, for the recalibrated IPSS-M and the elastic-net Cox model.
The diagonal marks perfect calibration.}
\label{fig5}
\end{figure}

\subsection*{External validation}
Table~\ref{table5} lists every submission to the test centre. The ensemble, specified as
the final model before submission, scored 0.7552, the highest of all submissions and within
0.0023 of its written prediction. The boosted model alone scored 0.7504, below the
elastic-net (0.7543) although it had beaten the elastic-net on 14 of 15 folds; a
subgroup analysis did not support the explanation that trees rank high-risk patients
poorly, since in the IPSS-R very-high-risk subgroup the boosted model ranked better than
the elastic-net in cross-validation (0.623 against 0.608, 11 of 15 folds). The IPSS-M scored
0.7255, 0.030 below the ensemble; no prediction was written down before this submission.
The fitted models' public scores were 0.022 to 0.033 higher than their cross-validated
means (Fig~\ref{fig6}), plausibly because the wider spread of risk at the test centre makes
patients easier to rank, but the IPSS-M gained only 0.005. Even the IPSS-M-style Cox model,
with the same variables as the IPSS-M but weights fitted to the training data, scored
higher at the test centre (0.7374), although it was lower in cross-validation (0.716).

The ensemble's lead over the IPSS-M was therefore six times larger at the external centre
than in cross-validation (0.030 against 0.005). Neither case mix nor missing data explained
this (Table~\ref{table6}). Validation folds resampled to the external centre's risk mix,
where 33\% of patients are IPSS-R very high risk against 11\% in training, gave the same
difference as before (0.005). Removing clinical values for an extra 7\% of patients lowered
the IPSS-M and the elastic-net almost equally (by 0.0026 and 0.0023). Within single training
centres, the difference ranged from $-$0.014 to 0.017 (standard deviation 0.011), and the
IPSS-M was ahead in three of eight centres, including the largest (Fig~\ref{fig7}). The
external centre's difference lay above this range.

\begin{table}[H]
\begin{adjustwidth}{-2.25in}{0in}
\centering
\caption{{\bf External validation on the unseen test centre.}}
\begin{tabular}{|l|c|c|c|}
\hline
{\bf Model} & {\bf Cross-validated C} & {\bf Written prediction} & {\bf Public score}\\ \thickhline
Elastic-net Cox & 0.7218 & -- & 0.7543 \\ \hline
IPSS-M-style Cox & 0.7159 & -- & 0.7374 \\ \hline
Ensemble (final model) & 0.7250 & 0.7575 & 0.7552 \\ \hline
Boosted Cox trees & 0.7244 & 0.7569 & 0.7504 \\ \hline
IPSS-M score (not fitted) & 0.7202 & -- & 0.7255 \\ \hline
\end{tabular}
\begin{flushleft} Public score: censoring-weighted concordance index computed by the
challenge platform on the public subset of the test centre. Written prediction: the
elastic-net's public score plus the model's cross-validated paired difference with the
elastic-net, written down before submission.
\end{flushleft}
\label{table5}
\end{adjustwidth}
\end{table}

\begin{figure}[H]
\plosfig{fig24_leaderboard.png}
\caption{{\bf Cross-validated against external concordance.}
Each point is one submitted model. Hollow markers show the predictions written down before
submission; dotted lines join them to the public scores.}
\label{fig6}
\end{figure}

\begin{table}[H]
\begin{adjustwidth}{-2.25in}{0in}
\centering
\caption{{\bf Why is the ensemble further ahead of the IPSS-M at the external centre?}}
\begin{tabular}{|l|c|l|}
\hline
{\bf Analysis} & {\bf Ensemble minus IPSS-M} & {\bf Detail}\\ \thickhline
Cross-validation, 15 folds & 0.005 & ensemble ahead on 11/15 folds \\ \hline
Folds resampled to the external case mix & 0.005 & ensemble ahead on 10/15 folds \\ \hline
Extra 7\% of clinical values removed & +0.0003$^{a}$ & change in the difference \\ \hline
Single training centres ($\geq$100 patients) & $-$0.014 to 0.017 & SD 0.011, 8 centres \\ \hline
External centre (public score) & 0.030 & 0.7552 minus 0.7255 \\ \hline
\end{tabular}
\begin{flushleft} Censoring-weighted concordance index truncated at 7 years. $^{a}$Elastic-net
minus IPSS-M, the change caused by removing the values.
\end{flushleft}
\label{table6}
\end{adjustwidth}
\end{table}

\begin{figure}[H]
\plosfig{fig29_ipssm_gap_by_centre.png}
\caption{{\bf The difference between the ensemble and the IPSS-M varies between centres.}
Ensemble minus IPSS-M concordance within each training centre with at least 100 patients
(out-of-fold predictions; number of patients in brackets) and at the external centre
(public score). The dashed line is the cross-validated difference.}
\label{fig7}
\end{figure}

\section*{Discussion}
In 3,173 patients from 23 centres, with validation on held-out centres and on an unseen
external centre, carefully built survival models reproduced the discrimination of the
published IPSS-M but did not clearly exceed it. The IPSS-M, used with its published
weights and no fitting, reached a cross-validated concordance of 0.720; the best fitted
model, an ensemble of an elastic-net and a boosted Cox model, reached 0.725, an advantage
whose confidence interval includes zero. The recalibrated IPSS-M was also slightly better
calibrated. At the external centre, however, the ensemble was 0.030 ahead of the IPSS-M.
Two readings fit this result. If the IPSS-M was developed partly on the training centres,
its cross-validated scores are optimistic, and a centre outside its development cohort
shows its true level; the IPSS-M-style model, refitted to the training data, also overtook
it there. Alternatively, the difference simply varies widely between hospitals, as it did
across training centres, and this centre lies at the far end. A single external centre
cannot separate the two.

Most of the predictive performance came from clinical knowledge rather than from the
algorithm. Moving from four generic variables to 24 variables chosen from the IPSS-R and
IPSS-M literature added 0.028; every later step, namely penalisation of 84 predictors,
boosted trees and the ensemble, added at most 0.006 more. This agrees with systematic
evidence that flexible algorithms rarely improve on well-specified regression in clinical
prediction~\cite{christodoulou2019}. The most important predictors, mutation burden,
cytogenetics and TP53, are those already built into the IPSS-M. The modest gain of the
boosted model, and the small trees preferred by its tuning, suggest that little structure
remains for trees to find once expert-designed variables such as the cytogenetic category
and TP53 multi-hit status are present.

Two practical points arose about validation across centres. First, harmonisation mattered:
the test centre sequenced a smaller gene panel and recorded karyotypes differently, and
mutation counts were the strongest signal only once restricted to genes sequenced at every
centre. Without this step a mutation count would have meant different things in training
and at the test centre. Second, between-centre variation was large compared with the
differences between models. The unfitted IPSS-R ranged from 0.624 to 0.729 across folds,
so comparisons of mean scores would have been uninformative; paired comparisons on the same
folds, with confidence intervals corrected for the overlap between folds~\cite{nadeau2003},
were needed to separate a consistent 0.003 gain from noise. The external centre also showed
that a model winning 14 of 15 internal folds can still lose at a new hospital, which is
the protection that averaging two different models provided here.

This study has limitations. The training centres very probably overlap with the cohort on
which the IPSS-M was developed~\cite{bernard2022}; if so, the IPSS-M's cross-validated
scores are partly in-sample and optimistic, and the external centre is its fairest test.
External validation rests on one centre, scored on a public subset of its patients, and
the platform reports no confidence interval for that score. The
data contain no age, treatment or disease subtype, all of which affect survival. Monocytes
and 60 genes could not be used because of differences between centres, and laboratory
calibration differences cannot be checked. The karyotype parser is a set of hand-written
rules rather than a validated nomenclature parser. The four-variable benchmark was
reconstructed from the challenge description. Finally, the boosted model's results differ
in the fourth decimal place between operating systems, although each machine reproduces
its own results exactly.

\section*{Conclusion}
For myeloid neoplasms, fitted survival models built on clinical, cytogenetic and molecular
data, and validated on hospitals they had not seen, matched the published IPSS-M in
cross-validation and exceeded it by 0.030 at one external centre, a difference not explained
by case mix or missing data but not yet confirmed beyond that centre. Claims that a new model improves prognosis should be tested
against the IPSS-M, with validation grouped by centre and paired comparisons, before they
are accepted.

\section*{Supporting information}

% Include only the SI item label in the paragraph heading. Use the \nameref{label} command to cite SI items in the text.
\paragraph*{S1 Table.}
\label{S1_Table}
{\bf TRIPOD+AI checklist.} Location in the manuscript of each item of the TRIPOD+AI
reporting guideline.

\paragraph*{S1 Appendix.}
\label{S1_Appendix}
{\bf Predictor definitions.} Karyotype parsing rules, IPSS-R points and cytogenetic
categories, IPSS-M coefficients and reference means, the shared gene panel and the pathway
gene lists.

\paragraph*{S1 Fig.}
\label{S1_Fig}
{\bf Cox models on each validation fold.} Concordance of each Cox model on each of the 15
centre-grouped folds, and the paired difference with the IPSS-R score.

\paragraph*{S2 Fig.}
\label{S2_Fig}
{\bf Tree models on each validation fold.} Concordance of the random survival forest and
the boosted Cox model on each fold, and the paired difference with the elastic-net.

\paragraph*{S3 Fig.}
\label{S3_Fig}
{\bf Time-varying effects.} Scaled Schoenfeld residuals plus the estimated coefficient for
NPM1 mutation and platelets, against the constant effect assumed by the Cox model.

\paragraph*{S4 Fig.}
\label{S4_Fig}
{\bf Time-dependent area under the curve.} Cumulative/dynamic area under the curve at 1, 3
and 5 years for each model (mean of 15 folds).

\paragraph*{S1 File.}
\label{S1_File}
{\bf Full technical report.} Every analysis stage, with all intermediate tables and
figures.

\section*{Acknowledgments}
I thank Qube Research and Technologies and Institut Gustave Roussy for organising the
challenge and making the data available, and the Challenge Data team at \'Ecole normale
sup\'erieure for hosting it.

\nolinenumbers

% Either type in your references using
% \begin{thebibliography}{}
% \bibitem{}
% Text
% \end{thebibliography}
%
% or
%
% Compile your BiBTeX database using our plos2015.bst
% style file and paste the contents of your .bbl file
% here. See http://journals.plos.org/plosone/s/latex for
% step-by-step instructions.
%
% Bibliography generated with plos2015.bst from references.bib and pasted here, as PLOS requires.
\begin{thebibliography}{10}

\bibitem{greenberg2012}
Greenberg PL, Tuechler H, Schanz J, Sanz G, Garcia-Manero G, Sol{\'e} F, et~al.
\newblock Revised international prognostic scoring system for myelodysplastic
  syndromes.
\newblock Blood. 2012;120(12):2454--2465.

\bibitem{schanz2012}
Schanz J, Tuechler H, Sol{\'e} F, Mallo M, Lu{\~n}o E, Cervera J, et~al.
\newblock New comprehensive cytogenetic scoring system for primary
  myelodysplastic syndromes ({MDS}) and oligoblastic acute myeloid leukemia
  after {MDS} derived from an international database merge.
\newblock J Clin Oncol. 2012;30(8):820--829.

\bibitem{bernard2022}
Bernard E, Tuechler H, Greenberg PL, Hasserjian RP, Arango~Ossa JE, Nannya Y,
  et~al.
\newblock Molecular international prognostic scoring system for myelodysplastic
  syndromes.
\newblock NEJM Evid. 2022;1(7):EVIDoa2200008.

\bibitem{ishwaran2008}
Ishwaran H, Kogalur UB, Blackstone EH, Lauer MS.
\newblock Random survival forests.
\newblock Ann Appl Stat. 2008;2(3):841--860.

\bibitem{friedman2001}
Friedman JH.
\newblock Greedy function approximation: a gradient boosting machine.
\newblock Ann Stat. 2001;29(5):1189--1232.

\bibitem{chen2016}
Chen T, Guestrin C.
\newblock {XGBoost}: a scalable tree boosting system.
\newblock In: Proceedings of the 22nd ACM SIGKDD International Conference on
  Knowledge Discovery and Data Mining. New York: ACM; 2016. p. 785--794.

\bibitem{christodoulou2019}
Christodoulou E, Ma J, Collins GS, Steyerberg EW, Verbakel JY, Van~Calster B.
\newblock A systematic review shows no performance benefit of machine learning
  over logistic regression for clinical prediction models.
\newblock J Clin Epidemiol. 2019;110:12--22.

\bibitem{austin2016}
Austin PC, van Klaveren D, Vergouwe Y, Nieboer D, Lee DS, Steyerberg EW.
\newblock Geographic and temporal validity of prediction models: different
  approaches were useful to examine model performance.
\newblock J Clin Epidemiol. 2016;79:76--85.

\bibitem{steyerberg2019}
Steyerberg EW.
\newblock Clinical prediction models: a practical approach to development,
  validation, and updating.
\newblock 2nd ed. Cham: Springer; 2019.

\bibitem{challenge162}
{Qube Research and Technologies}, {Institut Gustave Roussy}. Overall survival
  prediction for patients diagnosed with myeloid leukemia [data challenge 162];
  2025.
\newblock [cited 2026 Oct 24]. Available from:
  https://challengedata.ens.fr/challenges/162.

\bibitem{noren2016}
Noren DP, Long BL, Norel R, Rrhissorrakrai K, Hess K, Hu CW, et~al.
\newblock A crowdsourcing approach to developing and assessing prediction
  algorithms for {AML} prognosis.
\newblock PLoS Comput Biol. 2016;12(6):e1004890.

\bibitem{collins2024}
Collins GS, Moons KGM, Dhiman P, Riley RD, Beam AL, Van~Calster B, et~al.
\newblock {TRIPOD+AI} statement: updated guidance for reporting clinical
  prediction models that use regression or machine learning methods.
\newblock BMJ. 2024;385:e078378.

\bibitem{ipssmcalc}
{Papaemmanuil Lab}. ipssm: {IPSS-M} risk calculator, version 1.0.4 [software];
  2023.
\newblock [cited 2026 Oct 24]. Available from:
  https://www.npmjs.com/package/ipssm.

\bibitem{cox1972}
Cox DR.
\newblock Regression models and life-tables.
\newblock J R Stat Soc Series B Stat Methodol. 1972;34(2):187--220.

\bibitem{zou2005}
Zou H, Hastie T.
\newblock Regularization and variable selection via the elastic net.
\newblock J R Stat Soc Series B Stat Methodol. 2005;67(2):301--320.

\bibitem{simon2011}
Simon N, Friedman J, Hastie T, Tibshirani R.
\newblock Regularization paths for {Cox's} proportional hazards model via
  coordinate descent.
\newblock J Stat Softw. 2011;39(5):1--13.

\bibitem{uno2011}
Uno H, Cai T, Pencina MJ, D'Agostino RB, Wei LJ.
\newblock On the {C}-statistics for evaluating overall adequacy of risk
  prediction procedures with censored survival data.
\newblock Stat Med. 2011;30(10):1105--1117.

\bibitem{kaplan1958}
Kaplan EL, Meier P.
\newblock Nonparametric estimation from incomplete observations.
\newblock J Am Stat Assoc. 1958;53(282):457--481.

\bibitem{harrell1982}
Harrell FE, Califf RM, Pryor DB, Lee KL, Rosati RA.
\newblock Evaluating the yield of medical tests.
\newblock JAMA. 1982;247(18):2543--2546.

\bibitem{hung2010}
Hung H, Chiang CT.
\newblock Estimation methods for time-dependent {AUC} models with survival
  data.
\newblock Can J Stat. 2010;38(1):8--26.

\bibitem{graf1999}
Graf E, Schmoor C, Sauerbrei W, Schumacher M.
\newblock Assessment and comparison of prognostic classification schemes for
  survival data.
\newblock Stat Med. 1999;18(17--18):2529--2545.

\bibitem{vancalster2019}
Van~Calster B, McLernon DJ, van Smeden M, Wynants L, Steyerberg EW.
\newblock Calibration: the {Achilles} heel of predictive analytics.
\newblock BMC Med. 2019;17:230.

\bibitem{nadeau2003}
Nadeau C, Bengio Y.
\newblock Inference for the generalization error.
\newblock Mach Learn. 2003;52(3):239--281.

\bibitem{grambsch1994}
Grambsch PM, Therneau TM.
\newblock Proportional hazards tests and diagnostics based on weighted
  residuals.
\newblock Biometrika. 1994;81(3):515--526.

\bibitem{breiman2001}
Breiman L.
\newblock Random forests.
\newblock Mach Learn. 2001;45(1):5--32.

\bibitem{pedregosa2011}
Pedregosa F, Varoquaux G, Gramfort A, Michel V, Thirion B, Grisel O, et~al.
\newblock Scikit-learn: machine learning in {Python}.
\newblock J Mach Learn Res. 2011;12:2825--2830.

\bibitem{polsterl2020}
P{\"o}lsterl S.
\newblock scikit-survival: a library for time-to-event analysis built on top of
  scikit-learn.
\newblock J Mach Learn Res. 2020;21(212):1--6.

\bibitem{lifelines2019}
Davidson-Pilon C.
\newblock lifelines: survival analysis in {Python}.
\newblock J Open Source Softw. 2019;4(40):1317.

\end{thebibliography}

\end{document}
